\section{Background}

The emergence of agentic AI systems---autonomous software entities capable of reasoning, planning, spawning sub-agents, and executing multi-step workflows---has introduced a fundamental architectural challenge that existing frameworks are ill-equipped to address: how to manage the lifecycle, accountability, and provenance of agents that recursively spawn other agents. This section situates our work within five interconnected research threads: agent lifecycle management in multi-agent systems, accountability and provenance in distributed AI, fixed versus mutable delegation chain models, hierarchical task decomposition in AI planning, and the BX3 Framework as a substrate for immutable parent-chain architectures.

\subsection{Agent Lifecycle Management in Multi-Agent Systems}

Classical agent systems research established lifecycle as a core architectural concern, with seminal work defining agent states as spanning initialization, execution, communication, and termination phases \cite{ferber1999multi}. In modern agentic frameworks, lifecycle management has expanded to encompass dynamic spawning, resource budgeting, and hierarchical supervision. The \textbf{AgentSpawn} paradigm enables runtime collaboration where parent agents dynamically spawn specialized children based on task decomposition needs \cite{agentspawn2025}. Unlike training-time evolution systems, AgentSpawn operates at execution time, allowing orchestration layers to decide agent topology in response to workload characteristics rather than fixed configuration.

The architectural literature distinguishes between two dominant spawning patterns. The first is \textbf{declarative YAML-based configuration}, where sub-agent types are predefined with fixed name, system prompt, allowed tools, and context scope \cite{nibzard2025subagent}. The second is \textbf{dynamic runtime spawning}, where agents themselves determine when and how to spawn children based on real-time task demands \cite{microsoftconductor2025dynamic}. Dynamic spawning offers superior adaptability but introduces non-determinism in task graphs, making debugging and auditability more challenging. The Conductor framework proposes tool-based spawning with built-in safeguards including per-agent spawn limits (default 5) and workflow-level total agent caps (default 50), with spawned agents inheriting parent timeouts and token budgets. OpenClaw implements recursive spawning with configurable depth limits (default \texttt{maxSpawnDepth=1}, hard cap at 5 enforced by Zod schema validation), depth-aware model overrides (depth 0 $\rightarrow$ Opus, depth 1 $\rightarrow$ Sonnet, depth 2+ $\rightarrow$ Haiku), and orphan fallback mechanisms that surface orphaned sub-agent results to the root session when a parent is deleted \cite{openclaw2025recursive}.

Critical to lifecycle management is the problem of resource cleanup and state externalization. In Claude Code's architecture, each sub-agent invocation must balance providing sufficient context for the child to perform its task without incurring token waste on irrelevant data, while also enforcing permission boundaries strict enough for safety yet flexible enough for utility \cite{claudecode2025subagents}. The sub-agent lifecycle is implemented via two key components: \texttt{AgentTool.tsx} defining the model-facing interface and \texttt{runAgent.ts} implementing a multi-step execution lifecycle that manages environment creation, abort controller propagation, and result retrieval. The Multica framework enforces a \texttt{DEFAULT\_SUBAGENT\_TOOL\_DENY} policy that blocks sub-agents from spawning nested sub-agents, preventing uncontrolled recursion while preserving structured parent-child hierarchies \cite{multica2025sessionsspawn}.

The operational concern of \textbf{digital abandonment}---agents left running without active oversight---has emerged as a first-class lifecycle issue. Industry standards proposals call for clear ownership transfer protocols and sunset procedures with automatic shutdown mechanisms to prevent resources from being orphaned by completed or crashed parent agents \cite{astrasyncai2025abandonment}. This concern is especially acute in recursive spawning scenarios, where each additional depth level multiplies the potential for orphaned child sessions.

\subsection{Accountability and Provenance in Distributed AI Systems}

Accountability in multi-agent systems is fundamentally a provenance problem: every action taken by a spawned agent must be traceable to a human originator through an unbroken chain of delegation events. The concept of an \textbf{immutable delegation chain} has been formalized as the ordered record of every agent hop from the human who initiated a task down to the agent currently acting, where each link carries the agent's identity, its narrowed scope, its remaining budget, and the tools it is permitted to use \cite{agenticcontrolplane2025delegation}. DeepMind's treatment of intelligent AI delegation formalizes four accountability principles: (1) no humans required at inference time, (2) accountability is transitive across delegation hops, (3) provenance is immutable, and (4) in long chains (A $\rightarrow$ B $\rightarrow$ C), every downstream action remains attributable to the origin through the chain \cite{deepmind2025delegation}.

Blockchain-based approaches have been proposed to guarantee auditability in distributed AI systems. The core argument is that blockchain's immutable, decentralized, and transparent nature provides a powerful paradigm for establishing data provenance in AI auditing, guaranteeing that all auditing activities are traceable and verifiable in real-time \cite{blockchain2023accountability}. More specifically, Kite introduces cryptographically verifiable audit trails where every action creates an immutable chain of custody from user to agent to service \cite{cheng2025trustless}. This is particularly relevant for agentic systems, where the multi-hop delegation problem---where AI agents spawn other agents that call still other agents that access services---breaks down traditional OAuth-based authentication models that assume a single principal \cite{workos2025multihop}. Structured delegation tokens that encode both the user and agent identities along with permitted scopes at each hop offer a solution, enabling downstream services to validate the chain without contacting the authorization server at every hop \cite{scalekit2025delegated}. The IETF drafts addressing this gap propose nested claims encoding every level of authority from user to agent to service, ensuring every service in the chain can enforce policy and log actions properly.

The identity challenge compounds when agents are ephemeral and dynamic. Unlike static service accounts, agent identities must be established at spawn time, carry delegation scope, persist through execution, and be verifiable by downstream services---all without relying on pre-established credentials or shared secret material \cite{opal2025identity}. The concept of \textbf{zero-trust authorization} for multi-agent systems extends beyond the ``never trust, always verify'' principle to include delegation-specific controls: implementation of delegation limits before deploying multi-agent systems, cross-agent audit trails, and behavioral boundaries that prevent agents from accumulating permissions no single agent should hold \cite{gupta2025zerotrust}. A critical attack surface is the manipulation of delegation chains to accumulate permissions that no single agent in the chain individually possesses---a concern that requires validation of delegation chains at every hop rather than trust-on-first-use models. Cisco's identity framework for AI agents articulates that an agent's identity must inherently include its position in the delegation chain, its scope of authority, and the audit trail of its actions \cite{cisco2025identity}. The Okta treatment of Agentic RAG architecture extends this to real-time monitoring: logging every tool call and retrieval request with agent identity context, including which agent acted under whose delegation and what was accessed, to provide a complete audit trail across the lifecycle of agent identities that may themselves spawn other agents \cite{okta2025agenticrag}.

\subsection{Fixed vs. Mutable Delegation Chains}

The architectural decision between fixed and mutable delegation chains encodes a fundamental tradeoff between predictability and adaptability. \textbf{Fixed delegation chains} are established at system design time; the agent topology, permitted hops, and scope constraints are static and predetermined. This model appears in classical multi-agent systems where roles are pre-defined and agent interactions follow剧本-driven protocols \cite{ferber1999multi}. Fixed chains offer strong auditability---the full chain is known a priori---but cannot adapt to unexpected task complexity or runtime conditions that demand dynamic decomposition.

\textbf{Mutable delegation chains}, by contrast, allow agents to dynamically extend the chain by spawning sub-agents at runtime, with each spawn event potentially modifying the delegation scope, budget, and permitted toolset based on the parent agent's assessment of the child's task requirements \cite{agentspawn2025}. This model is dominant in modern agentic frameworks including CrewAI, LangGraph, and Microsoft's Conductor \cite{dhiwise2025frameworks}. The mutable model introduces what the industry calls the \textbf{delegation paradox}: as organizations transition from passive generative tools toward agentic systems, the need to delegate authority to agents conflicts with the need to maintain control and accountability over those delegations \cite{forbes2026delegation}. The paradox is acute in mutable chains because each spawn event potentially changes the scope and risk profile of the delegation without prior human review.

The IETF's work on OAuth 2.0 for multi-hop delegation highlights a core architectural flaw in existing authorization models: when Agent A spawns Agent B that calls Agent C that accesses Service D, the authorization server becomes a participant in every delegation decision, creating a bottleneck and a single point of failure \cite{workos2025multihop}. Structured delegation tokens that encode both the user and agent identities along with permitted scopes at each hop offer a solution, enabling downstream services to validate the chain without contacting the authorization server at every hop \cite{scalekit2025delegated}. However, this solution requires that each delegation token be immutable once issued---a property that is straightforward in fixed chains but non-trivial to guarantee in mutable chains where spawned agents may modify their scope before delegating further.

The BX3 Framework addresses this tension through a substrate-level commitment to immutable parent chains: the origin persists, scopes narrow at each hop (never expand), budget propagates but never increases, and every link in the chain is audited. This architectural commitment ensures that mutable spawning at the task level does not compromise the immutability of the accountability chain at the governance level.

\subsection{Hierarchical Task Decomposition in AI}

Hierarchical task decomposition is the process by which a complex task is recursively broken into simpler subtasks, forming a hierarchy that mirrors the agent spawning structure. The classical treatment is \textbf{Hierarchical Task Network (HTN) planning}, where a planner decomposes complex problems into smaller, structured subtasks through a formal grammar of tasks and methods \cite{geeksforgeeks2025htn}. HTN planning remains foundational in modern agent systems; classical planning systems such as STRIPS evolved into HTN, which are still essential for task decomposition in modern agents \cite{1password2025secure}. The key insight from HTN is that decomposition is not a one-time event but an iterative process: compound tasks are decomposed into subtasks, some of which may themselves be compound and require further decomposition, until primitive tasks with direct executable actions are reached.

The integration of HTN principles with large language models has opened new possibilities for runtime decomposition. The \textbf{ReCAP} (Recursive Context-Aware Reasoning and Planning) framework enables LLMs to perform long-horizon reasoning through three key mechanisms: plan-ahead decomposition, where the model generates a full subtask list, executes the first item, and refines the remainder; structured re-injection of parent plans to maintain consistent multi-level context during recursive return; and memory-efficient execution bounding active prompts so costs scale linearly with task depth \cite{recap2025neurips}. ReCAP's design maps directly onto recursive spawning: each decomposition event corresponds to a potential spawn, and the structured re-injection mechanism preserves the parent context in the child without requiring full context transfer.

The \textbf{ChatHTN} algorithm demonstrates a concrete integration of LLMs with HTN planning, where GPT is queried to generate task decompositions for compound tasks and the resulting subtask lists are recursively processed with a verifier task inserted at each level to validate decomposition correctness \cite{arxiv2505.11814}. This approach parallels the recursive spawning model presented here: the parent agent generates a decomposition plan, spawns specialized child agents to execute subtasks, and synthesizes results while maintaining the hierarchical context. The HTN planner's property that computational complexity is driven by task decomposition structure rather than problem size is directly relevant to recursive spawning depth management---a decomposition that spawns children with unbounded depth risks combinatorial explosion analogous to the planning complexity that motivated HTN research.

Task decomposition in LLM agents more broadly encompasses techniques including Chain-of-Thought reasoning \cite{chi2022chainofthought}, ReAct (Reason and Act) \cite{yao2022react}, and the Plan-and-Solve paradigm \cite{wang2023planandsolve}. These methods share the insight that breaking tasks into steps---guided by a decomposed subtask list---provides a foundation for complex multi-agent orchestration, as the structured execution path enables multiple specialized agents to operate on different subtask segments in parallel \cite{apxml2025taskdecomposition}. The hierarchical nature of this decomposition creates a natural mapping onto agent spawning topologies: a parent agent performing high-level task decomposition can treat each subtask as a candidate for delegation to a spawned child agent with specialized context and tools.

\subsection{The BX3 Framework as Substrate for Immutable Parent Chains}

The BX3 Framework provides a conceptual and structural substrate for agentic systems built around the principle that delegation chains must be immutable, accountable, and hierarchically structured from origin to leaf. The framework establishes a \textbf{chain-of-custody model} for AI agent governance: every spawned agent retains a reference to its parent, and every action taken by any agent in the chain is attributable through the chain to its origin. This model addresses the core failure mode identified in mutable delegation systems---where spawned agents modify their scope, accumulate permissions, or obscure the origin---and replaces it with a provenance-first architecture where immutability is a structural property, not a policy choice.

The BX3 framework's treatment of recursive spawning specifically addresses the problem that existing multi-agent systems treat spawning as an opaque execution event rather than a governed delegation act. In BX3, spawning is the primary governance event: it is the moment at which scope is narrowed, budget is propagated, tools are constrained, and identity is established for the child agent. The framework enforces that parent scopes cannot be exceeded by children (narrowing constraint), that budget propagates downward without augmentation (budget conservation), and that every spawn event produces an immutable chain link that persists through execution and audit (chain immutability).

Critically, BX3 decouples the task execution model from the accountability model: spawned agents can exhibit dynamic, adaptive behavior at the task level---spawning sub-agents, adapting strategies, invoking tools---while the accountability chain remains structurally immutable. This separation is essential for recursive spawning, where depth multiplication creates exponential complexity in the task graph but requires linear or constant overhead in the accountability structure. The framework's substrate provides this through what we term \textbf{parent-chain immutability}: the architectural commitment that the delegation chain from human origin to leaf agent is never mutable once established, even as the task execution hierarchy below any given node may grow and adapt dynamically.

This architectural property positions BX3 as a substrate for systems that require verifiable provenance without sacrificing the adaptive benefits of dynamic agent spawning. As agentic systems transition from single-agent tools to multi-agent ecosystems with recursive spawning capabilities, the framework provides the structural guarantees necessary for accountability, auditability, and governance at scale.\label{sec:recursive-spawning}